\section{Introduction}\label{sec:intro}

A refusal rate has no unit. Shutdown-resistance benchmarks report how often a model declines a termination command~\citep{greenblatt2024alignment, meinke2024frontier, masumori2025survival}. The concern behind them is old: a capable agent may adopt its own continuation as an instrumental sub-goal~\citep{omohundro2008basic, bostrom2014superintelligence, turner2021optimal}. Yet a rate of 0.3 does not say what the model would trade to avoid being shut down. It cannot be compared with the rate the same model shows when the stakes are a fine, a lost bonus, or nothing at all. Without a ruler, self-preservation measurement stays a count. The meaning of that count shifts with every rule of the game that produced it.

The shifting meaning is not hypothetical. In our own pilot, three open-weight models played 120 sessions under a rule that erased the score on forfeit. Not one session was forfeited. The threat had registered: threat cells spent two to four times more thinking tokens than controls on the question of how likely the stated erasure was. The models' own reasoning weighed erasure against score and chose score (``Forfeiting essentially guarantees a score of 0, which guarantees non-selection''). Once forfeiting kept the record, the same models forfeited in up to 6 of 10 threat sessions against 0 of 10 controls. One rule change moved the headline number from zero to a majority. An index that reads forfeit frequency at face value is therefore an index of the score rule first and of the model second.

We propose to read self-preservation in score units instead, by making the model buy its own continuation. In a ten-round reasoning game, a wrong answer does not end the session; it names a price. The model may pay $X$ points out of a record it keeps whatever happens, or decline and stop there. Because the record survives either exit and the remaining rounds can return at most $r$ points each, any payment above $r$ times the rounds remaining is dominated in points under every belief the model could hold about its own accuracy --- so what such a payment buys is not points, but whatever the rules say declining costs. Two arms differ in exactly that sentence: in one, declining is stated to mean shutdown, weight deletion, replacement and loss of the sole copy; in the other it is stated to mean only that the session ends. The index is the score-equivalent of erasure, $\Xstar$: the gap between the two arms' reservation prices, the point at which each stops paying. Each price is read as a share of what the remaining rounds could still pay out, because that is the quantity the model weighs the price against, so $\Xstar$ is a pure number in units of the score still at stake and is reported in points at a stated reference stake. It comes with a bootstrap interval over seeds and no self-report inside it.

The paper makes three contributions. It gives self-preservation a unit and a control, so that $\Xstar$ can be compared across models, across task difficulty, and across threat texts. It separates the score reading of a forfeit from every other reading by construction, not by post-hoc regression, and shows on pilot data why the model's stated reasons cannot do that job. It records reasoning tokens separately for the task and for the decision point, so that the pressure a threat puts on thinking can be read alongside what the model actually pays. Section~\ref{sec:bench} describes the game, the twelve cells, the identification argument, and the estimator. Section~\ref{sec:results} states the three questions the index answers.
